\chapter{Modules and Exact Sequences}
\label{ch:v09}

Modules generalize vector spaces by allowing scalars from an arbitrary ring. Exact sequences record how modules are assembled from submodules and quotients and become the basic grammar of homological algebra.

\section*{Learning goals}
After this chapter, the reader should be able to:
\begin{itemize}
\item verify exactness of sequences by computing kernels and images at each term;
\item construct short exact sequences from submodules and quotients;
\item apply the snake-lemma-style kernel/cokernel viewpoint in elementary diagrams where developed;
\item identify splitting maps and determine when a short exact sequence splits;
\item compute kernels, cokernels, images, and coimages of module homomorphisms;
\item distinguish left exact, right exact, and exact behavior of basic module functors;
\end{itemize}
\section*{Conceptual roadmap}
\volroadmap
  {presentation or universal property}
  {exactness/localization}
  {structural consequence}

\section{Modules}
An $A$-module is an abelian group with compatible scalar multiplication by $A$.

\section{Submodules and quotients}
Submodules are additive subgroups stable under scalars; quotient modules encode congruence modulo a submodule.

% BEGIN VOL05-EXPANSION V09-example-01
\begin{example}[Kernel and cokernel]\label{ex:v09-expand-01}
Multiplication by $6$ gives a short exact sequence $0\to\mathbb Z\to\mathbb Z\to\mathbb Z/6\to0$. The kernel is zero and the cokernel is the residue module.
\end{example}
% END VOL05-EXPANSION V09-example-01
\section{Module maps}
Kernels, images, and cokernels are modules and satisfy the familiar first isomorphism theorem.

\section{Exact sequences}
A sequence is exact when the image at each term equals the next kernel.

\section{Short exact sequences}
$0\to M'\to M\to M''\to0$ expresses $M'$ as a submodule and $M''$ as the corresponding quotient.

% BEGIN VOL05-EXPANSION V09-example-02
\begin{example}[Matrix exactness]\label{ex:v09-expand-02}
The inclusion $k^2\to k^3$ into the first two coordinates has cokernel $k$. Projection onto the third coordinate completes a short exact sequence.
\end{example}
% END VOL05-EXPANSION V09-example-02
\section{Splittings}
A short exact sequence splits when the middle module is isomorphic to a direct sum of its ends.

\section{Direct sums}
Direct sums collect finitely supported tuples and implement coproducts in the module category.

% BEGIN VOL05-EXPANSION V09-example-03
\begin{example}[Split sequence]\label{ex:v09-expand-03}
The sequence $0\to A\to A\oplus B\to B\to0$ splits via $b\mapsto(0,b)$, exhibiting the middle module as a direct sum.
\end{example}
% END VOL05-EXPANSION V09-example-03
\section{Diagram arguments}
Exact sequences organize kernel-image calculations and support diagram lemmas later in homological algebra.

\section{Examples}
Ideals, quotient modules, free modules, torsion abelian groups, and vector spaces all fit the same language.

\section{Core structural results}

\begin{theorem}[First isomorphism theorem for modules]\label{thm:v09-01}
For an $A$-linear map $f:M\to N$, $M/\ker f\cong\operatorname{im}f$.
\begin{proof}
Send $m+\ker f$ to $f(m)$. The same kernel argument as for groups and rings gives a well-defined bijective linear map onto the image.
\end{proof}
\end{theorem}

\begin{theorem}[Splitting criterion]\label{thm:v09-02}
A short exact sequence $0\to M'\xrightarrow{i}M\xrightarrow{p}M''\to0$ splits iff $p$ has a section, equivalently iff $i$ has a retraction.
\begin{proof}
A section $s$ gives $M=i(M')\oplus s(M'')$. Conversely a direct-sum decomposition supplies both projection and inclusion maps realizing the section and retraction.
\end{proof}
\end{theorem}

\begin{theorem}[Exactness of direct sums]\label{thm:v09-03}
Taking arbitrary direct sums of short exact sequences of modules preserves exactness.
\begin{proof}
All maps act componentwise. An element of a direct sum has finite support, so kernel and image membership can be checked component by component.
\end{proof}
\end{theorem}

\section{Worked examples}

\begin{example}[Modules diagnostic]\label{ex:v09-worked-01}
Audit Modules in four stages: hypotheses, decisive mechanism, argument, and limitation. For the decisive mechanism, track kernels and images at every arrow rather than reasoning only from the endpoints. An $A$-module is an abelian group with compatible scalar multiplication by $A$. The decisive algebraic mechanism is to track kernels and images at every arrow rather than reasoning only from the endpoints. The stated hypotheses determine which structural conclusion is valid.
\end{example}

\begin{example}[Submodules and quotients diagnostic]\label{ex:v09-worked-02}
For Submodules and quotients, begin by isolating the assumptions. Next, track kernels and images at every arrow rather than reasoning only from the endpoints. Conclude by stating precisely what has been proved and where the reasoning can break down. Submodules are additive subgroups stable under scalars; quotient modules encode congruence modulo a submodule. The decisive algebraic mechanism is to track kernels and images at every arrow rather than reasoning only from the endpoints. The stated hypotheses determine which structural conclusion is valid.
\end{example}

\begin{example}[Module maps diagnostic]\label{ex:v09-worked-03}
Treat Module maps as a short proof dossier. State the hypotheses explicitly, then track kernels and images at every arrow rather than reasoning only from the endpoints. Derive the conclusion and identify a failure mode if a required hypothesis is absent. Kernels, images, and cokernels are modules and satisfy the familiar first isomorphism theorem. The decisive algebraic mechanism is to track kernels and images at every arrow rather than reasoning only from the endpoints. The stated hypotheses determine which structural conclusion is valid.
\end{example}

\begin{example}[Exact sequences diagnostic]\label{ex:v09-worked-04}
Resolve Exact sequences by separating assumptions from inference. State the assumptions, track kernels and images at every arrow rather than reasoning only from the endpoints, give the conclusion, and note one condition under which that conclusion is no longer justified. A sequence is exact when the image at each term equals the next kernel. The decisive algebraic mechanism is to track kernels and images at every arrow rather than reasoning only from the endpoints. The stated hypotheses determine which structural conclusion is valid.
\end{example}

\section{Solved problems}
Each extended problem combines several results from the chapter in a longer argument. It is intended to make hypothesis selection, proof strategy, and verification explicit before the shorter exercises begin.

\begin{problem}[Modules diagnostic]\label{prob:v09-problem-01}
Audit Modules in four stages: hypotheses, decisive mechanism, argument, and limitation. For the decisive mechanism, track kernels and images at every arrow rather than reasoning only from the endpoints.
\end{problem}
\begin{solution}
An $A$-module is an abelian group with compatible scalar multiplication by $A$. The decisive algebraic mechanism is to track kernels and images at every arrow rather than reasoning only from the endpoints. The stated hypotheses determine which structural conclusion is valid.
\end{solution}

\begin{problem}[Submodules and quotients diagnostic]\label{prob:v09-problem-02}
For Submodules and quotients, begin by isolating the assumptions. Next, track kernels and images at every arrow rather than reasoning only from the endpoints. Conclude by stating precisely what has been proved and where the reasoning can break down.
\end{problem}
\begin{solution}
Submodules are additive subgroups stable under scalars; quotient modules encode congruence modulo a submodule. The decisive algebraic mechanism is to track kernels and images at every arrow rather than reasoning only from the endpoints. The stated hypotheses determine which structural conclusion is valid.
\end{solution}

\begin{problem}[Module maps diagnostic]\label{prob:v09-problem-03}
Treat Module maps as a short proof dossier. State the hypotheses explicitly, then track kernels and images at every arrow rather than reasoning only from the endpoints. Derive the conclusion and identify a failure mode if a required hypothesis is absent.
\end{problem}
\begin{solution}
Kernels, images, and cokernels are modules and satisfy the familiar first isomorphism theorem. The decisive algebraic mechanism is to track kernels and images at every arrow rather than reasoning only from the endpoints. The stated hypotheses determine which structural conclusion is valid.
\end{solution}

\begin{problem}[Exact sequences diagnostic]\label{prob:v09-problem-04}
Resolve Exact sequences by separating assumptions from inference. State the assumptions, track kernels and images at every arrow rather than reasoning only from the endpoints, give the conclusion, and note one condition under which that conclusion is no longer justified.
\end{problem}
\begin{solution}
A sequence is exact when the image at each term equals the next kernel. The decisive algebraic mechanism is to track kernels and images at every arrow rather than reasoning only from the endpoints. The stated hypotheses determine which structural conclusion is valid.
\end{solution}

\begin{problem}[Short exact sequences diagnostic]\label{prob:v09-problem-05}
Treat Short exact sequences as a short proof dossier. State the hypotheses explicitly, then track kernels and images at every arrow rather than reasoning only from the endpoints. Derive the conclusion and identify a failure mode if a required hypothesis is absent.
\end{problem}
\begin{solution}
$0\to M'\to M\to M''\to0$ expresses $M'$ as a submodule and $M''$ as the corresponding quotient. The decisive algebraic mechanism is to track kernels and images at every arrow rather than reasoning only from the endpoints. The stated hypotheses determine which structural conclusion is valid.
\end{solution}

\begin{problem}[Splittings diagnostic]\label{prob:v09-problem-06}
Treat Splittings as a short proof dossier. State the hypotheses explicitly, then track kernels and images at every arrow rather than reasoning only from the endpoints. Derive the conclusion and identify a failure mode if a required hypothesis is absent.
\end{problem}
\begin{solution}
A short exact sequence splits when the middle module is isomorphic to a direct sum of its ends. The decisive algebraic mechanism is to track kernels and images at every arrow rather than reasoning only from the endpoints. The stated hypotheses determine which structural conclusion is valid.
\end{solution}

\begin{problem}[Direct sums diagnostic]\label{prob:v09-problem-07}
Resolve Direct sums by separating assumptions from inference. State the assumptions, track kernels and images at every arrow rather than reasoning only from the endpoints, give the conclusion, and note one condition under which that conclusion is no longer justified.
\end{problem}
\begin{solution}
Direct sums collect finitely supported tuples and implement coproducts in the module category. The decisive algebraic mechanism is to track kernels and images at every arrow rather than reasoning only from the endpoints. The stated hypotheses determine which structural conclusion is valid.
\end{solution}

\begin{problem}[Diagram arguments diagnostic]\label{prob:v09-problem-08}
Treat Diagram arguments as a short proof dossier. State the hypotheses explicitly, then track kernels and images at every arrow rather than reasoning only from the endpoints. Derive the conclusion and identify a failure mode if a required hypothesis is absent.
\end{problem}
\begin{solution}
Exact sequences organize kernel-image calculations and support diagram lemmas later in homological algebra. The decisive algebraic mechanism is to track kernels and images at every arrow rather than reasoning only from the endpoints. The stated hypotheses determine which structural conclusion is valid.
\end{solution}

\begin{problem}[First isomorphism theorem for modules]\label{prob:v09-problem-09}
Prove the following structural statement and identify the hypothesis that makes the conclusion work: For an $A$-linear map $f:M\to N$, $M/\ker f\cong\operatorname{im}f$.
\end{problem}
\begin{solution}
Send $m+\ker f$ to $f(m)$. The same kernel argument as for groups and rings gives a well-defined bijective linear map onto the image. The proof hinges on the following algebraic step: track kernels and images at every arrow rather than reasoning only from the endpoints.
\end{solution}

\begin{problem}[Splitting criterion]\label{prob:v09-problem-10}
Prove the following structural statement and identify the hypothesis that makes the conclusion work: A short exact sequence $0\to M'\xrightarrow{i}M\xrightarrow{p}M''\to0$ splits iff $p$ has a section, equivalently iff $i$ has a retraction.
\end{problem}
\begin{solution}
A section $s$ gives $M=i(M')\oplus s(M'')$. Conversely a direct-sum decomposition supplies both projection and inclusion maps realizing the section and retraction. The proof hinges on the following algebraic step: track kernels and images at every arrow rather than reasoning only from the endpoints.
\end{solution}

\begin{problem}[Exactness of direct sums]\label{prob:v09-problem-11}
Prove the following structural statement and identify the hypothesis that makes the conclusion work: Taking arbitrary direct sums of short exact sequences of modules preserves exactness.
\end{problem}
\begin{solution}
All maps act componentwise. An element of a direct sum has finite support, so kernel and image membership can be checked component by component. The proof hinges on the following algebraic step: track kernels and images at every arrow rather than reasoning only from the endpoints.
\end{solution}

\begin{problem}[Chapter synthesis]\label{prob:v09-problem-12}
Connect the definitions and principal structural theorems of this chapter into one reusable commutative-algebra strategy.
\end{problem}
\begin{solution}
Modules generalize vector spaces by allowing scalars from an arbitrary ring. Exact sequences record how modules are assembled from submodules and quotients and become the basic grammar of homological algebra. A chapter-specific strategy is to track kernels and images at every arrow rather than reasoning only from the endpoints; consequently, exact sequences encode how one module is assembled from a submodule and a quotient.
\end{solution}

\section{Exercises with complete solutions}

\begin{exercise}\label{exr:v09-01}
For Modules, distinguish the formal setup from the decisive step. First record the hypotheses; then track kernels and images at every arrow rather than reasoning only from the endpoints. Explain the conclusion and a failure mode outside those hypotheses.
\end{exercise}
\begin{hint}
At each term of a sequence, compute image of the incoming map and kernel of the outgoing map; exactness means equality of those submodules.
\end{hint}
\begin{solution}
An $A$-module is an abelian group with compatible scalar multiplication by $A$. In this chapter the concrete consequence is that exact sequences encode how one module is assembled from a submodule and a quotient.
\end{solution}

\begin{exercise}\label{exr:v09-02}
For Submodules and quotients, distinguish the formal setup from the decisive step. First record the hypotheses; then track kernels and images at every arrow rather than reasoning only from the endpoints. Explain the conclusion and a failure mode outside those hypotheses.
\end{exercise}
\begin{hint}
For 0 -> N -> M -> Q -> 0, verify injectivity, surjectivity, and that Q is the quotient M/N up to the induced isomorphism.
\end{hint}
\begin{solution}
Submodules are additive subgroups stable under scalars; quotient modules encode congruence modulo a submodule. In this chapter the concrete consequence is that exact sequences encode how one module is assembled from a submodule and a quotient.
\end{solution}

\begin{exercise}\label{exr:v09-03}
Work through Module maps from the assumptions outward. Under the stated hypotheses, track kernels and images at every arrow rather than reasoning only from the endpoints. Give the resulting conclusion and one boundary case where the argument no longer applies.
\end{exercise}
\begin{hint}
A splitting of the surjection gives \(M\cong N\oplus Q\); construct the isomorphism explicitly from inclusion and section.
\end{hint}
\begin{solution}
Kernels, images, and cokernels are modules and satisfy the familiar first isomorphism theorem. In this chapter the concrete consequence is that exact sequences encode how one module is assembled from a submodule and a quotient.
\end{solution}

\begin{exercise}\label{exr:v09-04}
For Exact sequences, distinguish the formal setup from the decisive step. First record the hypotheses; then track kernels and images at every arrow rather than reasoning only from the endpoints. Explain the conclusion and a failure mode outside those hypotheses.
\end{exercise}
\begin{hint}
A splitting of the inclusion gives a retraction; use its kernel as the complementary summand.
\end{hint}
\begin{solution}
A sequence is exact when the image at each term equals the next kernel. In this chapter the concrete consequence is that exact sequences encode how one module is assembled from a submodule and a quotient.
\end{solution}

\begin{exercise}\label{exr:v09-05}
Determine which hypotheses are actually needed for Short exact sequences. Then track kernels and images at every arrow rather than reasoning only from the endpoints. State the conclusion and identify what can fail when one essential hypothesis is removed.
\end{exercise}
\begin{hint}
Compute cokernel as target modulo image, not merely as a set-theoretic complement.
\end{hint}
\begin{solution}
$0\to M'\to M\to M''\to0$ expresses $M'$ as a submodule and $M''$ as the corresponding quotient. In this chapter the concrete consequence is that exact sequences encode how one module is assembled from a submodule and a quotient.
\end{solution}

\begin{exercise}\label{exr:v09-06}
For Splittings, distinguish the formal setup from the decisive step. First record the hypotheses; then track kernels and images at every arrow rather than reasoning only from the endpoints. Explain the conclusion and a failure mode outside those hypotheses.
\end{exercise}
\begin{hint}
To test functor exactness, apply the functor to a short exact sequence and inspect where injectivity or surjectivity is lost.
\end{hint}
\begin{solution}
A short exact sequence splits when the middle module is isomorphic to a direct sum of its ends. In this chapter the concrete consequence is that exact sequences encode how one module is assembled from a submodule and a quotient.
\end{solution}

\begin{exercise}\label{exr:v09-07}
Use Direct sums as a hypothesis audit. List the assumptions, then track kernels and images at every arrow rather than reasoning only from the endpoints. State exactly what follows and what the argument would cease to control if an assumption were dropped.
\end{exercise}
\begin{hint}
For diagram arguments, chase one element at a time through commuting squares and write down where exactness is used.
\end{hint}
\begin{solution}
Direct sums collect finitely supported tuples and implement coproducts in the module category. In this chapter the concrete consequence is that exact sequences encode how one module is assembled from a submodule and a quotient.
\end{solution}

\begin{exercise}\label{exr:v09-08}
For Diagram arguments, distinguish the formal setup from the decisive step. First record the hypotheses; then track kernels and images at every arrow rather than reasoning only from the endpoints. Explain the conclusion and a failure mode outside those hypotheses.
\end{exercise}
\begin{hint}
Do not confuse exactness of a complex with splitting; a short exact sequence can be nonsplit.
\end{hint}
\begin{solution}
Exact sequences organize kernel-image calculations and support diagram lemmas later in homological algebra. In this chapter the concrete consequence is that exact sequences encode how one module is assembled from a submodule and a quotient.
\end{solution}

\section*{Chapter summary}
The chapter now contributes a worked-problem layer to the progression from rings and modules through localization, Noetherian methods, integral dependence, and derived functors.
% BEGIN VOL05-EXPANSION V09-exercises-01
\section{Graded supplementary exercises}
These exercises extend the chapter with a balanced set of computations, proofs, hypothesis tests, applications, and challenges.

\subsection*{Standard computations}

\begin{exercise}[Kernel]\label{exr:v09-expand-01}
Find the kernel of multiplication by $4$ on $\mathbb Z$.
\end{exercise}
\begin{hint}
Use kernels, cokernels, and exact sequences.
\end{hint}
\begin{solution}
It is zero.
\end{solution}

\begin{exercise}[Cokernel]\label{exr:v09-expand-02}
Find its cokernel.
\end{exercise}
\begin{hint}
Use kernels, cokernels, and exact sequences.
\end{hint}
\begin{solution}
It is $\mathbb Z/4$.
\end{solution}

\begin{exercise}[Injection]\label{exr:v09-expand-03}
Is multiplication by $2$ on $\mathbb Z$ injective?
\end{exercise}
\begin{hint}
Use kernels, cokernels, and exact sequences.
\end{hint}
\begin{solution}
Yes.
\end{solution}

\begin{exercise}[Matrix rank]\label{exr:v09-expand-04}
Find the rank of the projection $k^3\to k^2$.
\end{exercise}
\begin{hint}
Use kernels, cokernels, and exact sequences.
\end{hint}
\begin{solution}
The rank is $2$.
\end{solution}

\begin{exercise}[Section]\label{exr:v09-expand-05}
Give a section of $A\oplus B\to B$.
\end{exercise}
\begin{hint}
Use kernels, cokernels, and exact sequences.
\end{hint}
\begin{solution}
$b\mapsto(0,b)$.
\end{solution}

\subsection*{Proofs}

\begin{exercise}[First isomorphism theorem for modules proof]\label{exr:v09-expand-06}
Prove: For an $A$-linear map $f:M\to N$, $M/\ker f\cong\operatorname{im}f$.
\end{exercise}
\begin{hint}
Start from the chapter theorem hypotheses and identify the decisive algebraic map or containment.
\end{hint}
\begin{solution}
Send $m+\ker f$ to $f(m)$. The same kernel argument as for groups and rings gives a well-defined bijective linear map onto the image.
\end{solution}

\begin{exercise}[Splitting criterion proof]\label{exr:v09-expand-07}
Prove: A short exact sequence $0\to M'\xrightarrow{i}M\xrightarrow{p}M''\to0$ splits iff $p$ has a section, equivalently iff $i$ has a retraction.
\end{exercise}
\begin{hint}
Start from the chapter theorem hypotheses and identify the decisive algebraic map or containment.
\end{hint}
\begin{solution}
A section $s$ gives $M=i(M')\oplus s(M'')$. Conversely a direct-sum decomposition supplies both projection and inclusion maps realizing the section and retraction.
\end{solution}

\begin{exercise}[Exactness of direct sums proof]\label{exr:v09-expand-08}
Prove: Taking arbitrary direct sums of short exact sequences of modules preserves exactness.
\end{exercise}
\begin{hint}
Start from the chapter theorem hypotheses and identify the decisive algebraic map or containment.
\end{hint}
\begin{solution}
All maps act componentwise. An element of a direct sum has finite support, so kernel and image membership can be checked component by component.
\end{solution}

\begin{exercise}[Structural synthesis]\label{exr:v09-expand-09}
Prove a short corollary by combining First isomorphism theorem for modules with Splitting criterion. State every hypothesis you use.
\end{exercise}
\begin{hint}
Apply the first theorem to move to its structural description, then use the second theorem on that description.
\end{hint}
\begin{solution}
A valid synthesis first invokes First isomorphism theorem for modules and then applies Splitting criterion; the conclusion follows after checking the shared hypotheses.
\end{solution}

\subsection*{Counterexamples and hypothesis tests}

\begin{exercise}[Injective surjective]\label{exr:v09-expand-10}
Test the claim: Every injective module map is surjective.
\end{exercise}
\begin{hint}
Try a smallest standard ring or module from this chapter.
\end{hint}
\begin{solution}
False: multiplication by $2$ on $\mathbb Z$.
\end{solution}

\begin{exercise}[Surjective injective]\label{exr:v09-expand-11}
Test the claim: Every surjection is injective.
\end{exercise}
\begin{hint}
Try a smallest standard ring or module from this chapter.
\end{hint}
\begin{solution}
False: $\mathbb Z\to\mathbb Z/2$.
\end{solution}

\begin{exercise}[Zero composite]\label{exr:v09-expand-12}
Test the claim: $g\circ f=0$ implies exactness at the middle term.
\end{exercise}
\begin{hint}
Try a smallest standard ring or module from this chapter.
\end{hint}
\begin{solution}
False: it gives only image contained in kernel.
\end{solution}

\subsection*{Applications and investigations}

\begin{exercise}[Presentations]\label{exr:v09-expand-13}
How do exact sequences encode generators and relations?
\end{exercise}
\begin{hint}
Translate the situation into kernels, cokernels, and exact sequences.
\end{hint}
\begin{solution}
A free surjection supplies generators and its kernel records relations.
\end{solution}

\begin{exercise}[Failure measures]\label{exr:v09-expand-14}
What do kernel and cokernel measure?
\end{exercise}
\begin{hint}
Translate the situation into kernels, cokernels, and exact sequences.
\end{hint}
\begin{solution}
They measure failure of injectivity and surjectivity.
\end{solution}

\subsection*{Challenge problems}

\begin{exercise}[Local-global synthesis]\label{exr:v09-expand-15}
Connect the chapter ideas 'Modules' and 'Examples' in one rigorous argument.
\end{exercise}
\begin{hint}
State the relevant map, ideal, module, or universal property before drawing the conclusion.
\end{hint}
\begin{solution}
The argument begins with An $A$-module is an abelian group with compatible scalar multiplication by $A$. It then uses the later viewpoint: Ideals, quotient modules, free modules, torsion abelian groups, and vector spaces all fit the same language. The bridge is supplied by the structural theorems proved in the chapter.
\end{solution}

\begin{exercise}[Hypothesis audit]\label{exr:v09-expand-16}
Choose one theorem from the chapter and explain precisely what can fail if one key hypothesis is removed.
\end{exercise}
\begin{hint}
Use one of the explicit counterexamples in the graded set and compare it with the theorem statement.
\end{hint}
\begin{solution}
The theorem hypotheses are essential because the chapter's quotient, localization, exactness, or finiteness mechanism can fail outside them. A correct answer identifies the dropped hypothesis and exhibits a concrete failure.
\end{solution}

% END VOL05-EXPANSION V09-exercises-01
